% Editable report section: AI Tools and Development Methodology
\section{AI Tools and Development Methodology}
\label{sec:2}

I used ChatGPT/Codex with GPT-6.1 on ChatGPT Plus, together with Pi 0.86.0 using GPT-6.1 through the OpenAI Codex provider. AI support for the program covered specification extraction, architecture, implementation, tests and measurements. Work progressed from geometric primitives to the AVL beach line, event handling, cell construction and the interfaces, with checks at each stage.

The architecture followed the operations each component needed to support. The beach line needed ordered search, insertion and deletion; the event queue needed deterministic priority and a way to reject obsolete circle events. These requirements provided concrete criteria for assessing a proposed design before integrating its components.

\begin{table}[!htbp]
\centering
\caption{Checks used during AI-assisted implementation.}
\label{tab:1}
\begin{tabularx}{\linewidth}{@{}p{0.34\linewidth}X@{}}
\toprule
\textbf{Development activity} & \textbf{Verification used} \\
\midrule
Geometry and circle construction & Singular-limit, translation and equidistance tests \\
AVL beach line & Random insertion/deletion with structural invariants \\
Sweep and cells & Ownership, area and empty-circle checks \\
Interfaces and replay & CLI exports, callbacks, event stepping and pointer controls \\
Reference comparison & SciPy/Qhull and separate exact all-pairs clipping \\
Performance & Repeated construction timing and isolated memory processes \\
\bottomrule
\end{tabularx}
\end{table}

Figure~\ref{fig:ai-chatgpt} records a ChatGPT explanation of AVL order and generation-based circle-event rejection. The explanation was checked against the arc links, rotations and event guard. Pi also produced four regression tests covering stale predictions, duplicate circle delivery, successor-node identity during AVL deletion and deterministic queue ties. Those tests remain in \path{tests/test_pi_regressions.py}.

The second interaction, in Figure~\ref{fig:ai-pi}, examines a retained microscopic-cell failure. Its response connects the failure record with exact concurrence, circle identities and per-site assertions. The resulting checks verify that every distinct site retains a positive cell, including cases that area totals and sampled queries had missed. Both interactions concern specific code and test evidence.

Checks of generated code focused on geometric properties. Ownership validation first locates a query in the produced polygons and then compares the assigned sites with a separate nearest-site calculation. Empty-circle validation examines the emitted vertices independently of the drawing. The current reference comparison adds a second construction method and records where the reference library cannot resolve an input.

This process helped me connect the data structures to their geometric roles. A site can own several beach-line arcs, so its identifier cannot replace an arc's identity. A queued circle prediction must also match the current local configuration. Structural tests and retained counterexamples provided a way to assess these claims independently of the tool's explanations.
